% GENERATED FILE. Edit canonical lessons, then run npm run generate:books.
% canonical-source-hash: fnv1a64:36d3ae91fbbb722a
% canonical-lessons: HI-C249-sath, HI-C249-eksath, HI-C249-kampyutar, HI-C249-mobail, HI-C249-dastavez

\chapter{With, All together, A computer, A mobile phone, A document}
\label{ch:hi-with-all-together-a-computer-a-mobile-phone-a-document}
\hlchaptermodality{249}

\begin{chapteropening}
\textbf{By the end of this chapter:} \emph{I can say with, all together, a computer, a mobile phone and a document.}

Complete the last lesson of chapter 249: I can say with, all together, a computer, a mobile phone and a document.
\end{chapteropening}

\section[\texorpdfstring{\hi{साथ} (sāth)}{साथ (sāth)}]{\hi{साथ} (sāth) — with, together}

\label{lesson:HI-C249-sath}



\begin{quote}
Before the new one: say the Hindi for soon, then the Hindi for suddenly.
\end{quote}

\subsection*{You'll want to know: \hi{साथ}}
\textbf{\hi{साथ}} — \emph{sāth} — \textquotedblleft{}with, together\textquotedblright{}.

\begin{cousinweb}
\textbf{\hi{साथ}} (\emph{sāth}) is a tadbhava, a word inherited from Sanskrit and reshaped by sound change. It comes from \emph{sārthena}, \textquotedblleft{}in company with\textquotedblright{}, the instrumental case of \emph{sārtha}, \textquotedblleft{}caravan, company\textquotedblright{}. CDIAL lists Hindi \emph{sāth}, \textquotedblleft{}with\textquotedblright{}, under \emph{sārthena}, and the noun \emph{sāth}, \textquotedblleft{}company, train\textquotedblright{}, under \emph{sārtha} itself (CDIAL 13364).
\end{cousinweb}

\begin{grammarlens}[title={What you've built}]
One more word for this chapter.
\end{grammarlens}

\subsection*{Your turn}
\begin{itemize}
\raggedright
  \item \emph{Say it:} \emph{sāth}
  \item \emph{Say it:} \emph{sāth}, once more
  \item \emph{Say it:} say \emph{sāth}
  \item \emph{Recall:} say the Hindi for soon, then the Hindi for suddenly, then say \emph{sāth} again
\end{itemize}

\begin{tcolorbox}[breakable,colback=teal!4,colframe=teal!35!black,title={Before you move on}]
What does \hi{साथ} mean? (\textquotedblleft{}With, together\textquotedblright{}.) Say it once more.
\end{tcolorbox}
\section[\texorpdfstring{\hi{एक} \hi{साथ} (ek sāth)}{एक साथ (ek sāth)}]{\hi{एक} \hi{साथ} (ek sāth) — all together}

\label{lesson:HI-C249-eksath}



\begin{quote}
Before the new one: say the Hindi for suddenly, then the Hindi for with.
\end{quote}

\subsection*{You'll want to know: \hi{एक} \hi{साथ}}
\textbf{\hi{एक} \hi{साथ}} — \emph{ek sāth} — \textquotedblleft{}all together\textquotedblright{}.

\begin{cousinweb}
\textbf{\hi{एक} \hi{साथ}} (\emph{ek sāth}) is \emph{ek}, \textquotedblleft{}one\textquotedblright{}, from Sanskrit \emph{eka} (CDIAL 2462), and \emph{sāth}, \textquotedblleft{}with\textquotedblright{}, from \emph{sārthena}, \textquotedblleft{}in company with\textquotedblright{}, a case form of \emph{sārtha}, \textquotedblleft{}company\textquotedblright{} (CDIAL 13364).
\end{cousinweb}

\begin{grammarlens}[title={What you've built}]
One more word for this chapter.
\end{grammarlens}

\subsection*{Your turn}
\begin{itemize}
\raggedright
  \item \emph{Say it:} \emph{ek sāth}
  \item \emph{Say it:} \emph{ek sāth}, once more
  \item \emph{Say it:} say \emph{ek sāth}
  \item \emph{Recall:} say the Hindi for suddenly, then the Hindi for with, then say \emph{ek sāth} again
\end{itemize}

\begin{tcolorbox}[breakable,colback=teal!4,colframe=teal!35!black,title={Before you move on}]
What does \hi{एक} \hi{साथ} mean? (\textquotedblleft{}All together\textquotedblright{}.) Say it once more.
\end{tcolorbox}
\section[\texorpdfstring{\hi{कंप्यूटर} (kampyūṭar)}{कंप्यूटर (kampyūṭar)}]{\hi{कंप्यूटर} (kampyūṭar) — a computer}

\label{lesson:HI-C249-kampyutar}



\begin{quote}
Before the new one: say the Hindi for with, then the Hindi for all together.
\end{quote}

\subsection*{You'll want to know: \hi{कंप्यूटर}}
\textbf{\hi{कंप्यूटर}} — \emph{kampyūṭar} — \textquotedblleft{}a computer\textquotedblright{}.

\begin{cousinweb}
\textbf{\hi{कंप्यूटर}} (\emph{kampyūṭar}) is a loan from English \emph{computer}.
\end{cousinweb}

\begin{grammarlens}[title={What you've built}]
One more word for this chapter.
\end{grammarlens}

\subsection*{Your turn}
\begin{itemize}
\raggedright
  \item \emph{Say it:} \emph{kampyūṭar}
  \item \emph{Say it:} \emph{kampyūṭar}, once more
  \item \emph{Say it:} say \emph{kampyūṭar}
  \item \emph{Recall:} say the Hindi for with, then the Hindi for all together, then say \emph{kampyūṭar} again
\end{itemize}

\begin{tcolorbox}[breakable,colback=teal!4,colframe=teal!35!black,title={Before you move on}]
What does \hi{कंप्यूटर} mean? (\textquotedblleft{}A computer\textquotedblright{}.) Say it once more.
\end{tcolorbox}
\section[\texorpdfstring{\hi{मोबाइल} (mobāil)}{मोबाइल (mobāil)}]{\hi{मोबाइल} (mobāil) — a mobile phone}

\label{lesson:HI-C249-mobail}



\begin{quote}
Before the new one: say the Hindi for all together, then the Hindi for a computer.
\end{quote}

\subsection*{You'll want to know: \hi{मोबाइल}}
\textbf{\hi{मोबाइल}} — \emph{mobāil} — \textquotedblleft{}a mobile phone\textquotedblright{}.

\begin{cousinweb}
\textbf{\hi{मोबाइल}} (\emph{mobāil}) is a loan from English \emph{mobile}.
\end{cousinweb}

\begin{grammarlens}[title={What you've built}]
One more word for this chapter.
\end{grammarlens}

\subsection*{Your turn}
\begin{itemize}
\raggedright
  \item \emph{Say it:} \emph{mobāil}
  \item \emph{Say it:} \emph{mobāil}, once more
  \item \emph{Say it:} say \emph{mobāil}
  \item \emph{Recall:} say the Hindi for all together, then the Hindi for a computer, then say \emph{mobāil} again
\end{itemize}

\begin{tcolorbox}[breakable,colback=teal!4,colframe=teal!35!black,title={Before you move on}]
What does \hi{मोबाइल} mean? (\textquotedblleft{}A mobile phone\textquotedblright{}.) Say it once more.
\end{tcolorbox}
\section[\texorpdfstring{\hi{दस्तावेज़} (dastāvez)}{दस्तावेज़ (dastāvez)}]{\hi{दस्तावेज़} (dastāvez) — a document}

\label{lesson:HI-C249-dastavez}



\begin{quote}
Before the new one: say the Hindi for a computer, then the Hindi for a mobile phone.
\end{quote}

\subsection*{You'll want to know: \hi{दस्तावेज़}}
\textbf{\hi{दस्तावेज़}} — \emph{dastāvez} — \textquotedblleft{}a document\textquotedblright{}.

\begin{cousinweb}
\textbf{\hi{दस्तावेज़}} (\emph{dastāvez}) is a loan, usually traced to Persian \emph{dastāwēz}, \textquotedblleft{}a document, a deed\textquotedblright{}.
\end{cousinweb}

\begin{grammarlens}[title={What you've built}]
One more word for this chapter.
\end{grammarlens}

\subsection*{Your turn}
\begin{itemize}
\raggedright
  \item \emph{Say it:} \emph{dastāvez}
  \item \emph{Say it:} \emph{dastāvez}, once more
  \item \emph{Say it:} say \emph{dastāvez}
  \item \emph{Recall:} say the Hindi for a computer, then the Hindi for a mobile phone, then say \emph{dastāvez} again
\end{itemize}

\begin{tcolorbox}[breakable,colback=teal!4,colframe=teal!35!black,title={Before you move on}]
What does \hi{दस्तावेज़} mean? (\textquotedblleft{}A document\textquotedblright{}.) Say it once more.
\end{tcolorbox}
